\section{A genuine tangent partial-stable polynomial ideal is saturated by its actual coordinate variable}

All polynomials below are the actual multivariate polynomials over a
characteristic-zero field $k$; their variable set may be arbitrary.
The partial derivative is the literal existing polynomial derivation,
not an analytic differential or a formal-series replacement.

\begin{theorem}[Literal polynomial partial differentiation is locally nilpotent on every whole polynomial]
\label{hir:actual-polynomial-pderiv-locally-nilpotent}
\lean{Hironaka.HironakaPDerivStablePolynomialXSaturation.exists_polynomial_pderiv_pow_zero}
\leanok
For every actual polynomial $P$ and coordinate $i$, there is a finite
$N$ such that $D_i^N(P)=0$, where powers are the genuine endomorphism
powers of the original polynomial partial derivative linear map.
\end{theorem}
\begin{proof}
For a true monomial $aX^s$, induct on the actual exponent $s_i$.
If it is zero, the literal monomial derivative formula gives zero
at the first derivative. Otherwise that formula decreases the true
$i$ exponent by one, so the induction hypothesis on the resulting
monomial gives a finite vanishing iterate, and one more original
derivative proves the assertion. Every actual polynomial is a finite
sum of such monomials. For an addition of two already treated whole
polynomials, add their finite nilpotence bounds; higher powers still
vanish once the original bound vanishes. Linearity gives vanishing
of their sum. This finite polynomial induction gives a genuine
bound, not an input nilpotence or support certificate.
\end{proof}

\begin{theorem}[Exact successive original partial derivatives of coordinate multiplication]
\label{hir:actual-polynomial-pderiv-coordinate-multiple-iterate}
\lean{Hironaka.HironakaPDerivStablePolynomialXSaturation.polynomial_pderiv_pow_X_mul}
\leanok
For every $k\ge0$ and every whole polynomial $P$,
\[
 D_i^{k+1}(X_iP)=X_iD_i^{k+1}(P)+(k+1)D_i^k(P).
\]
All products, scalars, iterates and the coordinate variable in this
identity are the literal original polynomial operations.
\end{theorem}
\begin{proof}
For $k=0$, the original Leibniz rule and $D_iX_i=1$ give the formula.
Differentiate the formula at $k$. Actual linearity differentiates
the scalar term; Leibniz differentiates the coordinate multiple into
$D_i^{k+1}(P)+X_iD_i^{k+2}(P)$. The remaining term is
$(k+1)D_i^{k+1}(P)$, so combine these two genuine copies to obtain
$(k+2)D_i^{k+1}(P)$. Endomorphism power composition identifies each
successive derivative exactly and proves the induction step.
\end{proof}

\begin{theorem}[A genuine partial-stable original polynomial ideal is saturated by the same coordinate variable]
\label{hir:actual-pderiv-stable-polynomial-coordinate-saturation}
\uses{hir:actual-polynomial-pderiv-locally-nilpotent,hir:actual-polynomial-pderiv-coordinate-multiple-iterate}
\lean{Hironaka.HironakaPDerivStablePolynomialXSaturation.X_mul_mem_iff_of_pderiv_stable}
\leanok
Let $J$ be any actual polynomial ideal and suppose $D_iP\in J$ for
every whole $P\in J$. Then for every actual polynomial $P$,
$X_iP\in J$ if and only if $P\in J$.
No cylindrical-extension, normal-form, degree-bound or series
saturation certificate is a premise.
\end{theorem}
\begin{proof}
The reverse direction is ideal multiplication. For the forward
direction, genuine derivative stability implies that every original
iterate $D_i^{k+1}(X_iP)$ belongs to the same $J$. If
$D_i^{k+1}(P)\in J$, the preceding exact coordinate-multiple formula
and ideal subtraction therefore give $(k+1)D_i^k(P)\in J$.
Characteristic zero makes the field scalar $k+1$ nonzero; multiply
by its actual inverse to deduce $D_i^k(P)\in J$.

Choose the genuine finite nilpotence bound of $P$ from the first
theorem. The top derivative is zero and hence belongs to $J$.
Apply the proved step downward through that finite bound to reach
$D_i^0(P)=P\in J$. All memberships retain the same original ideal
and whole polynomial. This is polynomial coordinate saturation;
transfer to an entire formal initial cone or to genuine filtered
original analytic ideals requires their separate whole-ideal and
homogeneous-component dictionaries. No generic series saturation or
normal flatness is concluded here.
\end{proof}
